\subsection{正向极限的张量积}

设 \((\mathbf{A}_{\alpha}, \rho_{\beta\alpha})\) 是环的一个正向系统，且 \((\mathbf{E}_{\alpha}, f_{\beta\alpha})\)（相应地 \((\mathbf{F}_{\alpha}, g_{\beta\alpha})\)）是右（相应地，左）\(A_{\alpha}\)-模的一个正向系统。对 \(\alpha \leqslant \beta\)，存在一个Z-模同态

\[
f _ {\beta \alpha} \otimes g _ {\beta \alpha}: \mathrm{E} _ {\alpha} \otimes_ {\mathrm{A} _ {\alpha}} \mathrm{F} _ {\alpha} \rightarrow (\mathrm{E} _ {\beta}) _ {[ \mathrm{A} _ {\alpha} ]} \otimes_ {\mathrm{A} _ {\alpha}} (\mathrm{F} _ {\beta}) _ {[ \mathrm{A} _ {\alpha} ]}
\]

另一方面，存在一个典范的 Z-模同态

\[
\left(\mathrm{E} _ {\beta}\right) _ {[ \mathrm{A} _ {\alpha} ]} \otimes_ {\mathrm{A} _ {\alpha}} \left(\mathrm{F} _ {\beta}\right) _ {[ \mathrm{A} _ {\alpha} ]} \rightarrow \mathrm{E} _ {\beta} \otimes_ {\mathrm{A} _ {\beta}} \mathrm{F} _ {\beta}
\]

对应于环同态 \(\rho_{\beta \alpha}\) （§ 3，第 3 号，命题 2）；由此通过复合我们得到一个 Z-模同态

\[
h _ {\beta \alpha}: \mathrm{E} _ {\alpha} \otimes_ {\mathrm{A} _ {\alpha}} \mathrm{F} _ {\alpha} \rightarrow \mathrm{E} _ {\beta} \otimes_ {\mathrm{A} _ {\beta}} \mathrm{F} _ {\beta}
\]

它把张量积 \(x_{\alpha} \otimes y_{\alpha}\) 映为 \(f_{\beta\alpha}(x_{\alpha}) \otimes g_{\beta\alpha}(y_{\alpha})\)。显然

\[
\left(\mathrm{E} _ {\alpha} \otimes_ {\mathrm{A} _ {\alpha}} \mathrm{F} _ {\alpha}, h _ {\beta \alpha}\right)
\]

是由 \(\mathbf{Z}\)-模构成的一个正向系统。设 \(\mathbf{A} = \varinjlim \mathbf{A}_{\alpha}, \mathbf{E} = \varinjlim \mathbf{E}_{\alpha}, \mathbf{F} = \varinjlim \mathbf{F}_{\alpha}\)，并设 \(\rho_{\alpha}: \mathrm{A}_{\alpha} \to \mathrm{A}, f_{\alpha}: \mathrm{E}_{\alpha} \to \mathrm{E}, g_{\alpha}: \mathrm{F}_{\alpha} \to \mathrm{F}\) 为典范映射。如上所述，定义了一个Z-线性映射 \(\pi_{\alpha}: \mathrm{E}_{\alpha} \otimes_{\mathrm{A}_{\alpha}} \mathrm{F}_{\alpha} \to \mathrm{E} \otimes_{\mathrm{A}} \mathrm{F}\)，它把张量积 \(x_{\alpha} \otimes y_{\alpha}\) 映为 \(f_{\alpha}(x_{\alpha}) \otimes g_{\alpha}(y_{\alpha})\)，且显然这些映射构成一个正向系统。因此我们得到一个Z-线性映射

\[
\pi = \varinjlim \pi_ {\alpha}: \varinjlim (\mathrm{E} _ {\alpha} \otimes_ {\mathrm{A} \alpha} \mathrm{F} _ {\alpha}) \rightarrow \mathrm{E} \otimes_ {\mathrm{A}} \mathrm{F}.\tag{7}
\]

命题7. Z-线性映射(7)是双射的。

我们记 \(\mathrm{P} = \varinjlim (\mathrm{E}_{\alpha}\otimes_{\mathbf{A}_{\alpha}}\mathrm{F}_{\alpha})\)，并且对所有 \(\alpha \in \mathbf{I}\)，设 \(h_\alpha : \mathrm{E}_\alpha \otimes_{\mathbf{A}_\alpha}\mathrm{F}_\alpha \to \mathrm{P}\) 是典范映射。另一方面，对所有 \(\alpha \in \mathbf{I}\)，设

\[
t _ {\alpha}: \mathrm{E} _ {\alpha} \times \mathrm{F} _ {\alpha} \rightarrow \mathrm{E} _ {\alpha} \otimes_ {\mathrm{A} _ {\alpha}} \mathrm{F} _ {\alpha}
\]

是典范的Z-双线性映射；对于 \(\alpha \leqslant \beta\) ,

\[
t _ {\beta} \left(f _ {\beta \alpha} \left(x _ {\alpha}\right), g _ {\beta \alpha} \left(y _ {\alpha}\right)\right) = f _ {\beta \alpha} \left(x _ {\alpha}\right) \otimes g _ {\beta \alpha} \left(y _ {\alpha}\right) = h _ {\beta \alpha} \left(t _ {\alpha} \left(x _ {\alpha}, y _ {\alpha}\right)\right)
\]

，因此 \((t_{\alpha})\) 是映射的一个正向系统。把 \(\lim (\mathbf{E}_{\alpha}\times \mathbf{F}_{\alpha})\) 与 \(\mathbf{E}\times \mathbf{F}\) 典范地等同（集合论，III，§7，第7号，命题10），我们导出一个映射 \(t = \lim t_{\alpha}:\mathbf{E}\times \mathbf{F}\to \mathbf{P}\)，使得

\[
t \left(f _ {\alpha} (x _ {\alpha}), g _ {\alpha} (y _ {\alpha})\right) = h _ {\alpha} \left(t _ {\alpha} \left(x _ {\alpha}, y _ {\alpha}\right)\right) = h _ {\alpha} \left(x _ {\alpha} \otimes y _ {\alpha}\right).
\]

由集合论，III，§7，第5号，引理1，立即可以看出 \(t\) 是Z-双线性的；此外，对 \(x \in \mathbf{E}\), \(y \in \mathbf{F}\), \(\lambda \in \mathbf{A}\)，存在 \(\alpha \in I\)，使得 \(x = f_{\alpha}(x_{\alpha})\), \(y = g_{\alpha}(y_{\alpha})\), \(\lambda = \rho_{\alpha}(\lambda_{\alpha})\)，其中 \(\lambda_{\alpha} \in A_{\alpha}\), \(x_{\alpha} \in E_{\alpha}\), \(y_{\alpha} \in F_{\alpha}\)（集合论，III，§ 3，第 7 号，引理 1）；由此

\[
t (x \lambda , y) = h _ {\alpha} ((x _ {\alpha} \lambda_ {\alpha}) \otimes y _ {\alpha}) = h _ {\alpha} (x _ {\alpha} \otimes (\lambda_ {\alpha} y _ {\alpha})) = t (x, \lambda y).
\]

因此存在唯一的Z-线性映射 \(\pi':\mathbf{E} \otimes_{\mathbf{A}} \mathbf{F} \to \mathbf{P}\)，使得 \(\pi'(x \otimes y) = t(x, y)\)（§3，第1号，命题1）。此外，根据定义

\[
\pi^ {\prime} \left(\pi \left(h _ {\alpha} \left(x _ {\alpha} \otimes y _ {\alpha}\right)\right)\right) = \pi^ {\prime} \left(f _ {\alpha} \left(x _ {\alpha}\right) \otimes g _ {\alpha} \left(y _ {\alpha}\right)\right) = h _ {\alpha} \left(x _ {\alpha} \otimes y _ {\alpha}\right)
\]

\[
\pi (\pi^ {\prime} (f _ {x} (x _ {\alpha}) \otimes g _ {\alpha} (y _ {\alpha}))) = \pi (h _ {\alpha} (x _ {\alpha} \otimes y _ {\alpha})) = f _ {\alpha} (x _ {\alpha}) \otimes g _ {\alpha} (y _ {\alpha})
\]

且由于形如 \(f_{\alpha}(x_{\alpha}) \otimes g_{\alpha}(y_{\alpha})\)（相应地 \(h_{\alpha}(x_{\alpha} \otimes y_{\alpha})\)）的元素生成Z-模 E \(\otimes_{\mathbf{A}}\) F（相应地，P），故 \(\pi' \circ \pi\) 与 \(\pi \circ \pi'\) 是恒等映射。

粗略地说，命题7可表述为：张量积与正向极限交换，且通常（7）式的两边借助同构 \(\pi\) 加以等同。

推论1. 设 \((\mathbf{E}_{\alpha}', f_{\beta \alpha}')\) （相应地 \((\mathbf{F}_{\alpha}', g_{\alpha \beta}'))\) 是右（相应地，左）\(\mathbf{A}_{\alpha}\)-模的另一个正向系统；对所有 \(\alpha \in \mathbf{I}\)，设 \(u_{\alpha}: \mathbf{E}_{\alpha} \to \mathbf{E}_{\alpha}'\)（相应地 \(v_{\alpha}: \mathbf{F}_{\alpha} \to \mathbf{F}_{\alpha}'\)）是 \(\mathbf{A}_{\alpha}\)-线性映射，使得 \((u_{\alpha})\)（相应地 \((v_{\alpha})\)）是一个正向系统。则 \((u_{\alpha} \oplus v_{\alpha})\) 是 \(\mathbf{Z}\)-线性映射的一个正向系统，且图表

\[
\begin{array}{c c c} \varinjlim (\mathrm{E} _ {\alpha} \otimes_ {\mathrm{A} _ {\alpha}} \mathrm{F} _ {\alpha}) & \longrightarrow & (\varinjlim \mathrm{E} _ {\alpha}) \otimes_ {\mathrm{A}} (\varinjlim \mathrm{F} _ {\alpha}) \\ \varinjlim (u _ {\alpha} \otimes v _ {\alpha}) \Bigg \downarrow & & \Bigg \downarrow (\varinjlim u _ {\alpha}) \otimes (\varinjlim v _ {\alpha}) \\ \varinjlim (\mathrm{E} _ {\alpha} ^ {\prime} \otimes_ {\mathrm{A} _ {\alpha}} \mathrm{F} _ {\alpha} ^ {\prime}) & \longrightarrow & (\lim \mathrm{E} _ {\alpha} ^ {\prime}) \otimes_ {\mathrm{A}} (\lim \mathrm{F} _ {\alpha} ^ {\prime}) \end{array}\tag{8}
\]

是交换的。

验证是直接的。

设 \((\mathbf{A}_{\alpha}^{\prime}, \rho_{\beta\alpha}^{\prime})\) 是环的另一个正向系统，并假设每个 \(E_{\alpha}\) 是一个 \((\mathbf{A}_{\alpha}^{\prime}, \mathbf{A}_{\alpha})\)-双模，且诸 \(f_{\beta\alpha}\) 是 \((\mathbf{A}_{\alpha}^{\prime}, \mathbf{A}_{\alpha})\)-线性的（对 \(\alpha \leqslant \beta\)）。则若记 \(A^{\prime} = \varinjlim A_{\alpha}^{\prime}\)，由推论1，同构（7）对两侧的左 \(A^{\prime}\)-模结构是线性的。这可立即推广到任意多模。

特别地，若诸 \(A_{\alpha}\) 交换，则 \(A = \varinjlim A_{\alpha}\) 交换，且同构（7）是A-模同构。

推论2. 设 \((\mathbf{E}_{\alpha}, f_{\beta \alpha})\) 是右 \(\mathbf{A}_{\alpha}\)-模的一个正向系统，并令 \(\mathbf{E}_{\alpha}' = \mathbf{E}_{\alpha} \otimes_{\mathbf{A}_{\alpha}} \mathbf{A}\) 为把标量环扩张到 \(\mathbf{A} = \lim_{\mathbf{A}_{\alpha}} \mathbf{A}_{\alpha}\)（借助典范同态 \(\rho_{\alpha}: \mathbf{A}_{\alpha} \to \mathbf{A}\)）而得的A-模。则 \((\mathbf{E}_{\alpha}', f_{\beta \alpha} \otimes 1_{\mathbf{A}})\) 是右A-模的一个正向系统，其正向极限典范同构于 \(\lim_{\mathbf{E}_{\alpha}} \mathbf{E}_{\alpha}\)。

只需在命题7中取 \(F_{\alpha}\) 为视为 \((\mathbf{A}_{\alpha}, \mathbf{A})\)-双模（借助 \(\rho_{\alpha}\)）的环A。

推论3. 设A是一个环， \((\mathrm{E}_{\alpha}, f_{\beta\alpha})\) 是右A-模的一个正向系统，且F是一个左A-模。则Z-模 \(\lim_{\longrightarrow} (\mathrm{E}_{\alpha} \otimes_{\mathrm{A}} \mathrm{F})\) 和 \((\lim_{\longrightarrow} \mathrm{E}_{\alpha}) \otimes_{\mathrm{A}} \mathrm{F}\) 典范同构。

只需在命题7中取 \(A_{\alpha} = A\) 和 \(F_{\alpha} = F\)（对所有 \(\alpha \in I\)）。

特别地，若 \(\rho: \mathbf{A} \to \mathbf{B}\) 是环同态，则 \(\lim_{\rho^* (\mathbf{E}_\alpha)} \rho^*(\mathbf{E}_\alpha)\) 和 \(\rho^*(\lim_{\rho \to \infty} \mathbf{E}_\alpha)\) 典范同构。

推论4. 设M为右A-模，N为左A-模， \((x_{i})_{1\leqslant i\leqslant n}\) 为M的元素族， \((y_{i})_{1\leqslant i\leqslant n}\) 为N的元素族，使得 \(\sum_{i}(x_{i}\otimes y_{i}) = 0\)（在 M \(\otimes_{\mathbf{A}}\mathbf{N}\) 中）。则存在M（相应地N）的一个有限生成子模 \(\mathbf{M}_1\)（相应地 \(\mathbf{N}_1\)），它包含诸 \(x_{i}\)（相应地，诸 \(y_{i}\)），且使得 \(\sum_{i}(x_{i}\otimes y_{i}) = 0\)（在 \(\mathbf{M}_1\otimes_{\mathbf{A}}\mathbf{N}_1\) 中）。

M（相应地N）被典范地等同于其包含诸 \(x_{i}\)（相应地，诸 \(y_{i}\)）的有限生成子模构成的右有向族的正向极限，于是只需应用集合论，III，§ 7，第 5 号，引理 1。

